\section{Qualitative midpoint arguments in infinite height}\label{sec:qualitative-infinite}

We work in the coordinate predual $B_\infty$ of
Section~\ref{binf:definition}, with the distinction from the full dual
specified in Remark~\ref{presentation:infinite-height}. The averaged
modulus is defined in \eqref{eq:averaged-modulus}. Three different
reductions lead to a vanishing-deficit contradiction: adding a
supporting kernel, normalizing the displacement twice, and
normalizing the two endpoints separately.

The common recovery step is Theorem~\ref{stab:weighted}. To make its
inputs explicit, let two finite laws of represented atoms have means
$w_\pm$, and choose a head test $p=P_Hp$ with $\rho(p)\le1$. Put
\[
 \delta_\pm=1-w_\pm(p),\qquad
 \delta=\frac{\delta_++\delta_-}{2},\qquad
 \eta=\left\|Q_H\frac{w_++w_-}{2}\right\|.
\]
Use the toolkit's squared-test slack $D$ and squared coefficient error
$E$ for these laws and this test. The pointwise budget and
Theorem~\ref{stab:weighted}, with weights $1/2$, give
\begin{align}
 \mathbb E_\pm(D+E)&\le2\delta_\pm,
       \label{qual:conditional-budget}\\
 \|Q_Hw_\pm\|&\le
 2\sqrt{(\eta+\sqrt{2\delta})^2+4\delta}
       +\sqrt{2\delta_\pm}.
       \label{qual:conditional-recovery}
\end{align}
Thus $\delta_\pm\to0$ and $\eta\to0$ force both conditional tails
to vanish, even when the heads, tests, and finite laws vary. The first
two reductions below give an exactly zero pooled tail. Separate
endpoint normalization gives a residual tail whose norm must also be
shown to tend to zero. The conditional budgets record what each
normalization preserves in the atom representations.

\subsection{A qualitative contradiction with a supporting kernel}
\label{mold:qualitative}
Suppose that \(\widehat\delta_{B_\infty}(t)=0\) for some
\(0<t<1\).  For each positive integer \(n\), choose a unit center
\(v\in B_\infty\) whose centerwise supremum in the definition of this
modulus is less than \(1/n\), and a finitely supported unit vector \(x\) with
\(\|x-v\|<1/n\).  Let \(A\) be a finite initial set supporting \(x\), and
choose \(f\in B_\infty^*\) with \(\|f\|=f(v)=1\).  In
\(\ker P_A\cap\ker f\), choose \(y\) of norm at least one such that
\[
 \frac{\|v+ty\|+\|v-ty\|}{2}<1+2/n.
\]
The supporting-functional condition makes each endpoint norm at least one;
hence both are less than \(1+4/n\).  Approximate \(ty\) by a finitely
supported \(u\in\ker P_A\) within \(\min\{t/2,1/n\}\), as allowed by
Lemma~\ref{binf:projections}.  Then
\[
 \|u\|\ge t/2,\qquad
 \max\{\|x+u\|,\|x-u\|\}\le1+6/n.
\]

Choose \(r\ge1\) dominating these endpoint norms, with \(r\to1\), and
use Lemma~\ref{binf:finite-hull} to represent \((x\pm u)/r\) exactly
by finite convex combinations of atoms in a finite ancestral set containing
their supports. Choose a head test \(p=P_Ap\) with
\(\rho(p)=x(p)=1\). For the means \(w_\pm=(x\pm u)/r\),
\[
 \delta_\pm=1-1/r\longrightarrow0,\qquad
 Q_A\frac{w_++w_-}{2}=Q_A(x/r)=0.
\]
In particular, each law has
\(\mathbb E_\pm(D+E)\le2(1-1/r)\to0\).
Equation~\eqref{qual:conditional-recovery} now gives
\(\|u\|/r=\|Q_Aw_+\|\to0\), contradicting
\(\|u\|\ge t/2\). Thus the averaged modulus is positive.

\subsection{A qualitative route with two normalizations}
\label{mold:qualitative-normalized}
A second reduction uses the radial monotonicity of the symmetric average.
Suppose that \(\widehat\delta_{B_\infty}(t)=0\) at a fixed
\(0<t<1\).  Choose unit centers in \(B_\infty\) whose
centerwise suprema tend to zero; approximate each by a finitely supported
unit vector \(x\) with error tending to zero, and let \(H\) be a finite
initial set containing its support.  In \(\ker P_H\), the infimum provides
a vector \(y\) of norm at least one for which the endpoint average at
\(x\), with displacement \(ty\), is
\[
 \frac{\|x+ty\|+\|x-ty\|}{2}\le1+o(1).
\]
For a fixed direction, the symmetric average
\[
 s\longmapsto\frac{\|x+sy\|+\|x-sy\|}{2}
\]
is even and convex, hence nondecreasing for \(s\ge0\).  Replacing \(y\)
by \(y/\|y\|\) therefore does not increase that average.  Next approximate
this unit vector in \(\ker P_H\) by finite vectors, using
Lemma~\ref{binf:projections}, and normalize the approximants again.
Their errors can be made to tend to zero, so we obtain
finite unit \(y\in\ker P_H\) with the same \(1+o(1)\) bound.

Both endpoint norms are at least one by \(P_H\), and therefore both tend
to one.  Choose \(r\) strictly larger than their maximum, with \(r\to1\).
The finite vectors \((x\pm ty)/r\) have norm strictly less than one, so
the strict gauge assertion in Lemma~\ref{binf:finite-hull} gives exact
finite convex representations. Choose a head test \(p=P_Hp\) with
\(\rho(p)=x(p)=1\). Each mean \(w_\pm=(x\pm ty)/r\) has evaluation
\(1/r\), so its own coefficient budget is
\(\mathbb E_\pm(D+E)\le2(1-1/r)\to0\).
The pooled mean is \(x/r\), whose tail vanishes. Thus
\(\delta_\pm=1-1/r\to0\) and \(\eta=0\) in
\eqref{qual:conditional-recovery}. It follows that
\(\|Q_Hw_\pm\|\to0\), whereas these norms equal \(t/r\to t>0\).
This contradiction proves positivity.

\subsection{A sequential argument with separate endpoint normalizations}
\label{molc:sequential-section}

We now prove the averaged midpoint property of $B_\infty$
qualitatively, normalizing the two endpoints separately.

We use the averaged modulus defined in \eqref{eq:averaged-modulus},
with the inner infimum taken over directions of norm at least one.
Its inner expression is nonnegative by the triangle inequality.

\begin{proposition}\label{molc:sequential-positive}
For each fixed $0<t<1$, the space $B_\infty$ satisfies
$\widehat\delta_{B_\infty}(t)>0$.
\end{proposition}

\begin{proof}
Suppose instead that $\widehat\delta_{B_\infty}(t)=0$. There are unit vectors $x_n$
for which the supremum over finite-codimensional subspaces is less than
$1/n$. Approximate $x_n$ within $1/n$ by a finitely supported unit vector
$a_n$, and let $A_n$ be a finite nonempty initial set containing its
support. Taking the subspace $\ker P_{A_n}$ in the definition, choose
$y_n$ there with $\|y_n\|\ge1$ and
\[
 \frac{\|x_n+ty_n\|+\|x_n-ty_n\|}{2}\le1+2/n.
\]
By Lemma~\ref{binf:projections}, choose a finite vector
$z_n\in\ker P_{A_n}$ within $1/n$ of $ty_n$. It follows that
\[
 \|z_n\|\ge t-1/n,\qquad
 \frac{\|a_n+z_n\|+\|a_n-z_n\|}{2}\le1+4/n.
\]
Put $r_n^\pm=\|a_n\pm z_n\|$. Projection onto the head shows that
$r_n^\pm\ge1$, so both sequences tend to $1$. Moreover,
\begin{equation}\label{molc:sequential-tail-bounded}
 \|z_n\|\le\frac{r_n^++r_n^-}{2}\le1+4/n.
\end{equation}
This boundedness is needed when the two normalizations are compared.

For each $n$, take a finite initial set $B_n\supseteq A_n$ containing
both supports. Choose a head test $p_n=P_{A_n}p_n$ with
$\rho(p_n)=a_n(p_n)=1$. Lemma~\ref{binf:finite-hull} gives two finite laws
of represented molecules $G=\sum_i\lambda_i\mathbf1_{S_i}$ satisfying
\[
 \mathbb E_{n,\pm}G=(a_n\pm z_n)/r_n^\pm.
\]
Write $w_{n,\pm}$ for these means. Their deficits on $p_n$ are
$\delta_{n,\pm}=1-1/r_n^\pm\to0$. With $D,E$ computed using $p_n$,
each law retains the budget
\begin{equation}\label{molc:sequential-deficits}
 \mathbb E_{n,\pm}(D+E)
 \le2\left(1-(\mathbb E_{n,\pm}G)(p_n)\right)
 =2(1-1/r_n^\pm)\longrightarrow0.
\end{equation}
Write $Q_n=I-P_{A_n}$ and let $\mathbb E_n$ denote the equally
weighted mixture. Its tail is
\begin{equation}\label{molc:unequal-normalization-tail}
 \mathbb E_n Q_nG
 =\frac12\left(\frac1{r_n^+}-\frac1{r_n^-}\right)z_n.
\end{equation}
It need not vanish exactly. Its norm $\eta_n$ tends to zero because
$r_n^\pm\to1$ and \eqref{molc:sequential-tail-bounded} bounds
$\|z_n\|$. Apply \eqref{qual:conditional-recovery} with the head
$A_n$, the test $p_n$, the deficits $\delta_{n,\pm}$, and this
$\eta_n$. It gives $\|Q_nw_{n,\pm}\|\to0$. These norms are
$\|z_n\|/r_n^\pm$, which cannot tend to zero because
$\|z_n\|\ge t-1/n$ and $r_n^\pm\to1$. This contradiction proves the
proposition.
\end{proof}

The obstruction to an equivalent asymptotically uniformly convex norm
is Proposition~\ref{binf:obstruction}, an instance of the general
criterion in Proposition~\ref{foundation:bounded-tree}.
